\subsection{极大秩子群}

回忆（第六章，§1，第 7 节）：R = R(G, T) 的子集 P 称为闭的，如果 \((P + P) \cap R \subset P\)；称为对称的，如果 P = -P。

命题 12. 设 H 是 G 的包含 T 的连通闭子群全体所成的集合，按包含关系排序。映射 \(\mathrm{H} \mapsto \mathrm{R}(\mathrm{H}, \mathrm{T})\) 是从 H 到 \(\mathrm{R}(\mathrm{G}, \mathrm{T})\) 的对称闭子集全体所成的集合（按包含关系排序）上的一个递增双射。

若 \(H \in H\)，则 \(\mathrm{L}(\mathrm{H})_{(\mathbf{C})}\) 是 \(t_{C}\) 与诸 \(g^{\alpha}\)（\(\alpha \in \mathrm{R}(\mathrm{H}, \mathrm{T})\)）的直和；由于它是 \(g_{C}\) 中的一个约化子代数，故 R 的子集 \(\mathrm{R}(\mathrm{H}, \mathrm{T})\) 满足所述条件（第八章，§3，第 1 节，引理 2 与命题 2）。反之，若 P 是满足这些条件的 R 的子集，则 \(t_{C} \oplus \sum_{\alpha \in P} g^{\alpha}\) 是 \(g_{C}\) 的一个子代数（同前处），它在 R 上是有理的（第 3 节），因而形如 \(h_{(C)}\)，其中 h 是 g 的子代数。设 \(\mathrm{H}(\mathrm{P})\) 是由 h 定义的 G 的整子群；它是闭的（§2，第 4 节，注 1）。我们立即验证：映射 \(\mathrm{H} \mapsto \mathrm{R}(\mathrm{H}, \mathrm{T})\) 与 \(\mathrm{P} \mapsto \mathrm{H}(\mathrm{P})\) 都是递增的且互为逆映射。

推论 1. G 中包含 T 的闭子群只有有限多个。

设 H 是这样一个子群；则 \(H_{0} \in H\)，且 H 有限。此外，H 是 \(\mathrm{N}_{\mathrm{G}}(\mathrm{H}_{0})\) 的包含 \(H_{0}\) 的子群，而 \(\mathrm{N}_{\mathrm{G}}(\mathrm{H}_{0})/\mathrm{H}_{0}\) 有限（\(\S2\)，第 4 节，命题 4 与注 2）。

推论 2. 设 H 是 G 的包含 T 的连通闭子群，并设 \(W_{\mathrm{G}}^{\mathrm{H}}(\mathrm{T})\) 是 \(W_{\mathrm{G}}(\mathrm{T})\) 中 R 的子集 \(\mathrm{R}(\mathrm{H},\mathrm{T})\) 的稳定子群。群 \(N_{\mathrm{G}}(\mathrm{H})/\mathrm{H}\) 同构于商群 \(W_{\mathrm{G}}^{\mathrm{H}}(\mathrm{T})/W_{\mathrm{H}}(\mathrm{T})\)。

实际上，把 §2 第 5 节命题 7 应用于 \(\mathrm{N}_{\mathrm{G}}(\mathrm{H})\)，可知 \(\mathrm{N}_{\mathrm{G}}(\mathrm{H})/\mathrm{H}\) 同构于 \(\mathrm{W}_{\mathrm{N}(\mathrm{H})}(\mathrm{T})/\mathrm{W}_{\mathrm{H}}(\mathrm{T})\)，其中 \(\mathrm{W}_{\mathrm{N}(\mathrm{H})}(\mathrm{T})\) 是 \(\mathrm{W}_{\mathrm{G}}(\mathrm{T})\) 中其代表在 \(\mathrm{N}_{\mathrm{G}}(\mathrm{T})\) 中且正规化 H 的元素全体所成的集合。设 \(n \in \mathrm{N}_{\mathrm{G}}(\mathrm{T})\)，并设 w 是它在 \(\mathrm{W}_{\mathrm{G}}(\mathrm{T})\) 中的类。由第三章 §9 第 4 节命题 11，n 正规化 H 当且仅当 \((\operatorname{Ad} n)(\operatorname{L}(\operatorname{H})) = \operatorname{L}(\operatorname{H})\)；由第 3 节命题 5，这也意味着 R 的子集 \(\operatorname{R}(\operatorname{H}, \operatorname{T})\) 在 w 下稳定，由此得推论。

注 1. 群 \(W_{\mathrm{G}}^{\mathrm{H}}(\mathrm{T})\) 也是 \(W_{\mathrm{G}}(\mathrm{T})\) 中 T 的子群 \(\mathrm{C}(\mathrm{H})\) 的稳定子群：这由第 4 节命题 8 得出。

命题 13. 设 H 是 G 的一个极大秩连通闭子群，C 是它的中心。则 C 包含 G 的中心，且 H 是 C 的中心化子的单位元连通分支。

设 S 是 H 的一个极大环面。由于 G 的中心含于 S，故它含于 C。令 \(L = Z(C)_{0}\)；这是 G 的一个包含 H 的连通闭子群，故为极大秩，且其中心等于 C。用 \(R_{H}\) 与 \(R_{L}\) 分别表示 H 与 L 关于 S 的根系；则 \(R_{H} \subset R_{L} \subset R(G, S)\)。由于 \(C(H) = C(L)\)，命题 8（第 4 节）蕴含等式 \(Q(R_{H}) = Q(R_{L})\)；但 \(Q(R_{H}) \cap R_{L} = R_{H}\)（第六章，§1，第 7 节，命题 23），故 \(R_{H} = R_{L}\)，从而 H = L（命题 12）。

注 2. 称子群 \(\mathrm{C}\)（\(\mathrm{G}\) 中）为根子群，如果存在极大环面 \(\mathrm{S}\)（\(\mathrm{G}\) 中）与子集 \(\mathrm{P}\)（\(\mathrm{R}(\mathrm{G},\mathrm{S})\) 中），使得 \(\mathrm{C} = \bigcap_{\alpha \in \mathrm{P}}\mathrm{Ker}\alpha\)。由此

并结合命题 13 与第 5 节引理 2 可知，映射 \(H \mapsto C(H)\) 诱导出从极大秩连通闭子群全体所成的集合到 G 的根子群全体所成的集合的双射。其逆双射是映射 \(C \mapsto Z(C)_{0}\)。

推论. \(g \in \mathbf{G}\) 中满足 \(\mathrm{T} \cap g\mathrm{T}g^{-1} \neq \mathrm{C}(\mathrm{G})\) 者所成的集合，是 \(\mathbf{G}\) 中异于 \(\mathbf{G}\) 的有限多个闭解析子流形的并。

实际上，令 \(A_{g} = T \cap gTg^{-1}\)；我们有 \(T \subset Z(A_{g})\) 与 \(gTg^{-1} \subset Z(A_{g})\)。于是存在 \(x \in Z(A_{g})\) 使得 \(xTx^{-1} = gTg^{-1}\)（§2，第 2 节，定理 2），这蕴含 \(g \in Z(A_{g}).N_{G}(T)\)。用 A 表示 G 中包含 T 且异于 G 的闭子群所成的有限集合（推论 1），并令 \(X = \bigcup_{H \in A} H.N_{G}(T)\)；这是 G 的有限多个闭子流形的并，且异于 G。若 \(A_{g} \neq C(G)\)，则 \(Z(A_{g}) \in A\)，故 g 属于 X。反之，若 \(g \in H.N_{G}(T)\)（其中 \(H \in A\)），则 \(A_{g}\) 包含 \(C(H)\)，故 \(A_{g} \neq C(G)\)（命题 13）。

命题 14. 设 X 是 T 的一个子集，并设 \(R_{X}\) 是根 \(\alpha \in \mathrm{R}(\mathrm{G}, \mathrm{T})\) 中满足 \(\alpha(\mathrm{X}) = \{1\}\) 者所成的集合。群 \(\mathrm{Z}_{\mathrm{G}}(\mathrm{X}) / \mathrm{Z}_{\mathrm{G}}(\mathrm{X})_{0}\) 同构于 \(\mathrm{W}_{\mathrm{G}}(\mathrm{T})\) 中稳定 X 的子群关于由诸反射 \(s_{\alpha}\)（\(\alpha \in R_{X}\)）所生成的子群所得的商。

令 \(H = Z_{\mathrm{G}}(\mathrm{X})\)；由于 \(\mathrm{L}(\mathrm{H})_{(\mathbf{C})}\) 是 \(g_{C}\) 中被 \(\mathrm{Ad}(\mathrm{X})\) 固定的点所成的集合，它是 \(t_{C}\) 与诸 \(g^{\alpha}\)（满足 \(\alpha(\mathrm{X}) = \{1\}\) 者）的和。因此，\(\mathrm{R}(\mathrm{H}_{0}, \mathrm{T}) = \mathrm{R}_{\mathrm{X}}\)，故 \(\mathrm{W}_{\mathrm{H}_{0}}(\mathrm{T})\) 由诸反射 \(s_{\alpha}\)（\(\alpha \in R_{X}\)）生成。此时只需应用 §2 第 5 节命题 7。

我们将在下文（\(\S5\)，第 3 节，定理 1）看到：若 G 单连通且 X 仅由一点组成，则中心化子 \(Z(X)\) 是连通的。

